\documentclass[11pt]{article}

% --- packages -------------------------------------------------------------
\usepackage[utf8]{inputenc}
\usepackage[T1]{fontenc}
\usepackage{amsmath,amssymb,amsthm}
\usepackage{booktabs}
\usepackage{graphicx}
\usepackage{hyperref}
\usepackage{geometry}
\geometry{margin=1in}
\usepackage{microtype}

\newcommand{\E}{\mathrm{E}}
\newcommand{\LE}{\mathrm{LE}}
\newcommand{\Wfast}{W_{\mathrm{fast}}}
\newcommand{\Wslow}{W_{\mathrm{slow}}}
\newcommand{\T}{T_{80}}

\title{\bfseries Developmental Learning Architecture:\\[2pt]
Fast/Slow Memory, Difficulty Robustness, and the Search for Learning-Rule Development}
\author{Anonymous\\ \texttt{https://github.com/JayCRL/DLA}}
\date{Draft v0.1}

\begin{document}
\maketitle

\begin{abstract}
We study a family of neural architectures in which the learning rule itself is a developmental object.
Our Developmental Learning Architecture (DLA) keeps a standard Transformer/MLP skeleton unchanged but attaches per-parameter fast weights, plasticity traces, and sleep consolidation, so that knowledge and the ability to acquire knowledge are separated into slow and fast state.
We report four main findings.
First, DLA's fast/slow separation reproduces the adaptation quality of tuned AdamW fine-tuning while keeping consolidated slow memory essentially immune to forgetting in sequential language-domain adaptation (5 seeds).
Second, DLA is robust to a progressively harder curriculum: its normalized learning efficiency does not collapse when task difficulty increases.
Third, a causal dissection shows that learning history strongly changes future learning performance on an unseen domain, but the 9-dimensional learning-rule parameters (tempos) are \emph{not} the carrier of this effect.
Fourth, body-state cross-injection localises the carrier to fast weights: transplanting fast weights from a hard$\to$easy learner into an easy$\to$hard body changes the future adaptation \emph{trajectory} rather than the initial perplexity.
However, with 10 seeds this fast-weight transfer effect is directionally consistent but not statistically conclusive.
We frame this paper as an honest mechanistic study: learning history shapes the learner, the effect is carried by accumulated fast-weight state, and current meta-updates of scalar learning-rule parameters do not yet produce stable improvements in future learning.
\end{abstract}

\section{Introduction}\label{sec:intro}
Large models are typically treated as fixed functions that are trained and then deployed.
This paper asks a different question: \emph{can the way a model learns change as a result of what it has learned?}
Drawing on ideas of developmental plasticity and fast/slow memory, we build a minimal architecture with:
\begin{itemize}
  \item \textbf{Slow weights} --- consolidated long-term knowledge;
  \item \textbf{Fast weights} --- a current learning trace;
  \item \textbf{Plasticity} $P$ --- a per-parameter state controlling how much each weight can change;
  \item \textbf{Consolidation} $Q$ --- an eligibility trace used during sleep.
\end{itemize}
We deliberately do \emph{not} modify the Transformer skeleton.
The developmental state lives beside the standard weights and is updated by an online learning rule, with periodic sleep consolidation.

The central empirical question is whether learning history can shape the \emph{future learner itself}, not only the knowledge it stores.
We operationalise this by comparing learners that experienced the same three domains in different orders (easy$\to$hard versus hard$\to$easy) and then measuring how they adapt to an unseen fourth domain.

\section{Related Work}\label{sec:related}
\begin{itemize}
  \item \textbf{Fast weights and memory.} Classical and recent fast-weight mechanisms \cite{ba2016fast,hinton1987using} motivate separating fast and slow stores.
  \item \textbf{Metaplasticity and continual learning.} Stability--plasticity trade-offs and synaptic tagging are studied in neuroscience \cite{abraham2005memory} and machine learning \cite{zenke2017continual}.
  \item \textbf{Learning to learn / meta-learning.} ANML \cite{beaulieu2020learning} and differentiable plasticity \cite{miconi2019differentiable} make selective plasticity learnable.
  \item \textbf{Learned optimizers / learning rules.} Learned optimizers \cite{andrychowicz2016learning} and differentiable learning rules make the update rule itself trainable.
  \item \textbf{This paper.} Existing meta-learning asks whether a good rule can be learned; we ask whether, within a single lifetime, the learning rule and its associated body state continue to develop, and which state actually carries development forward.
\end{itemize}

\section{Method}\label{sec:method}
\subsection{DLA on an MLP core}
For small-scale mechanism experiments we use a two-layer MLP.
Every connection has a slow weight $\Wslow$, a fast weight $\Wfast$, plasticity $P$, and an eligibility trace $Q$.
The effective weight is
\begin{equation}
  W_{\mathrm{eff}} = \Wslow + \mathrm{softplus}(P) \odot \Wfast .
\end{equation}
The experience update (``wake'') uses a per-connection Learning Rule Network $F_\phi$ that receives local and cognitive signals and outputs coefficients for a Hebbian basis and a teaching basis ($\delta x^\top$).
Sleep consolidates $Q$ into $\Wslow$ and decays fast traces.
Meta-training unrolls mini-lifetimes of two tasks with a retention loss, so the stability--plasticity trade-off is part of the objective.

\subsection{DLA on a standard GPT}
For Transformer experiments we wrap each linear layer and embedding of a standard nanoGPT model (6 layers, 8 heads, 256 dimensions, 7280 Chinese characters) with per-matrix states $\Wfast, P, Q$ and Adam moments.
The backbone code is not changed; each forward pass uses $W_{\mathrm{eff}}$ and each backward pass supplies the teaching signal $g = \mathrm{d}\mathcal L / \mathrm{d} W_{\mathrm{eff}}$.

The nine scalar tempos
\begin{equation}
  \phi = \{\eta_{\mathrm{fast}}, \eta_{\mathrm{plast}}, \lambda_{\mathrm{decay}}, \kappa_{\mathrm{stab}}, \alpha_q, \text{consolidation rates}\}
\end{equation}
start as DNA priors.
In a ``meta'' variant they can be updated by one meta-gradient step on future-task losses.

\subsection{Causal dissection}
We decompose the developmental state (``body'') into slow weights, fast weights, plasticity $P$, and eligibility $Q$.
Cross-injection swaps one component between two individuals with different histories while holding the rest fixed.
We then measure future adaptation on a held-out domain $D$ under identical data, update rule, budget, and evaluation.

\section{Experimental Setup}\label{sec:setup}
\begin{itemize}
  \item Backbone: 6.59M Chinese character-level GPT, pretrained on $\sim$9.3M tokens of simplified Chinese Wikipedia.
  \item Curriculum domains: general Wikipedia, SFT-style QA, science Wikipedia.
  \item Held-out domain $D$: a science-Wikipedia slice never used in the curriculum.
  \item Each individual trains each domain with the same budget ($40$--$100$ Adam-style steps of $32\times128$ tokens depending on stage).
  \item Metrics: raw relative gain; normalized learning efficiency $\LE_t = G_t(T)/G_t^{\max}$; $\T_t$ (steps to reach $80\%$ of the individual's own best gain); backward transfer/forgetting; plasticity and state traces.
\end{itemize}

\section{Results}\label{sec:results}
\subsection{Fast/slow separation protects consolidated memory (5 seeds)}\label{sec:fastslow}
\begin{table}[h]\centering\small
\begin{tabular}{lccc}
\toprule
Metric & AdamW (tuned) & DLA fast & DLA slow \\
\midrule
B-domain adaptation gain & $+13.0\% \pm 2.4$ & $+14.0\% \pm 2.5$ & $+0.1\%$ \\
A forgetting after B & $+10.1\% \pm 2.2$ & $+7.3\% \pm 1.2$ & $\mathbf{-0.4\% \pm 0.2}$ \\
\bottomrule
\end{tabular}
\caption{Stage 4 formal, 5 seeds. DLA matches tuned AdamW on adaptation while the consolidated slow memory remains essentially stable.}
\end{table}

\subsection{Difficulty robustness (5 seeds)}\label{sec:robust}
Under normalized learning efficiency,
\begin{table}[h]\centering\small
\begin{tabular}{lccc}
\toprule
Arm & $\LE_1$ & $\LE_2$ & $\LE_3$ \\
\midrule
AdamW & 0.63 & 0.87 & 0.25 \\
DLA fast & 0.79 & 0.70 & 0.69 \\
DLA slow & 0 & 0 & 0 \\
\bottomrule
\end{tabular}
\caption{Stage 5, 5 seeds. DLA does not collapse when task difficulty increases.}
\end{table}

\subsection{Scalar learning-rule parameters are not the carrier (5.5b/c, 5 seeds)}\label{sec:notphi}
\begin{table}[h]\centering\small
\begin{tabular}{lcc}
\toprule
Body / $\phi$ & $\LE_D$ mean \\
\midrule
EH/EH & 0.522 \\
EH/HE & 0.148 \\
HE/EH & 0.898 \\
HE/HE & 0.900 \\
\bottomrule
\end{tabular}
\caption{2$\times$2 cross-injection (body $\times$ phi) on unseen $D$. The body dominates; swapping the 9 tempos has little and inconsistent effect.}
\end{table}

\subsection{Fast weights carry part of the history effect (5.5d/5.5e, 10 seeds)}\label{sec:wfast}
Body decomposition (3 seeds) showed:
\begin{itemize}
  \item EH body $+$ HE fast weights $\Rightarrow$ $\LE_D$ increases from $0.33$ to $0.91$;
  \item HE body $+$ EH fast weights $\Rightarrow$ $\LE_D$ drops from $0.96$ to $0.00$.
\end{itemize}

Extended trajectory experiment (10 seeds):
\begin{table}[h]\centering\small
\begin{tabular}{lccc}
\toprule
Combo & gain@40 & sign contrast \\
\midrule
EH/EH & $-0.001$ & --- \\
HE/HE & $+0.014$ & --- \\
EH body $+$ HE fast & $+0.005$ & $+0.006$ vs.\ EH/EH (6/10 positive) \\
HE body $+$ EH fast & $-0.007$ & $-0.021$ vs.\ HE/HE (8/10 negative) \\
\bottomrule
\end{tabular}
\caption{Fast-weight transfer changes future adaptation slope more than initial perplexity. With 10 seeds the positive contrast is not statistically significant (sign-test $p\approx0.38$); the destructive contrast is borderline ($p\approx0.055$).}
\end{table}

\section{Discussion}\label{sec:discussion}
\textbf{What we now believe.}
\begin{enumerate}
  \item A developmental architecture can decouple fast adaptation from slow memory preservation.
  \item A hard$\to$easy curriculum produces a learner that adapts faster on an unseen domain than an easy$\to$hard learner; this history effect is reproducible.
  \item The effect is mediated by accumulated body state, especially fast weights, not by scalar meta-learned tempos.
  \item Fast-weight transfer changes the \emph{dynamics} of future adaptation, not merely the initial perplexity.
\end{enumerate}

\textbf{What we do not yet claim.}
\begin{itemize}
  \item That meta-updating scalar learning rules yields stable lifelong improvements in transfer (5.5b/5.5c negative/nuanced).
  \item That fast-weight transfer alone is a sufficient causal carrier (5.5e underpowered).
  \item That this scales beyond a 6.59M model.
\end{itemize}

\textbf{Limitations.}
Single small backbone; single language; synthetic domain slices; no standard continual-learning baselines (EWC, replay, LoRA) yet; 5--10 seeds on noisy PPL metrics; sleep consolidation writes very little to slow memory in the current implementation.

\section{Conclusion}\label{sec:conclusion}
We present an honest mechanistic account of a developmental learning architecture.
Positive contributions are (a) fast/slow memory separation, (b) difficulty robustness, and (c) evidence that learning history shapes the future learner and that this shaping is carried in part by fast-weight state.
Negative contributions are equally important: scalar learning-rule parameters are not a stable carrier of development, and a clean causal proof that fast weights change future learning dynamics still requires more power and better baselines.

\paragraph{Reproducibility.}
Code is at \url{https://github.com/JayCRL/DLA}.
Result JSONs and logs live at \texttt{\~/llm-lab/dla-v0.2/results/} and \texttt{/logs/}.
Key scripts: \texttt{stage4\_formal.py}, \texttt{stage5\_progressive\_curriculum.py}, \texttt{stage55*}, \texttt{stage55d\_body\_decomposition.py}, \texttt{stage55e\_wfast\_p0.py}.
Full per-stage tables: \texttt{docs/experiments.md}.

\begin{thebibliography}{99}
\bibitem{abraham2005memory} W.~Abraham and A.~Robins. Memory retention: the synaptic stability versus plasticity dilemma. \emph{Trends in Neurosciences}, 2005.
\bibitem{zenke2017continual} F.~Zenke, B.~Poole, and S.~Ganguli. Continual learning through synaptic intelligence. \emph{ICML}, 2017.
\bibitem{ba2016fast} J.~Ba, G.~Hinton, V.~Mnih, J.~Leibo, and C.~Ionescu. Using fast weights to attend to the recent past. \emph{NeurIPS}, 2016.
\bibitem{beaulieu2020learning} S.~Beaulieu, L.~Frantar, E.~Mironov, Y.~Chen, and R.~Savarese. Learning to continually learn. \emph{ECAI}, 2020.
\bibitem{miconi2019differentiable} T.~Miconi, J.~Clune, and K.~Stanley. Differentiable plasticity: training plastic neural networks with backpropagation. \emph{ICLR}, 2019.
\bibitem{hinton1987using} G.~Hinton and D.~Plaut. Using fast weights to deblur old memories. \emph{CogSci}, 1987.
\bibitem{andrychowicz2016learning} M.~Andrychowicz et al. Learning to learn by gradient descent by gradient descent. \emph{NeurIPS}, 2016.
\end{thebibliography}

\end{document}
